\chapter{Análisis del problema}\label{cap:analisis}

\section{Introducción}
En este capítulo detallamos los requisitos fundamentales que el sistema logístico debe satisfacer. El análisis se ha abordado desde una perspectiva integral, considerando tanto la robustez necesaria en la persistencia y el cálculo de datos en el backend  como la agilidad requerida en la interfaz de usuario para una gestión operativa real.

\section{Requisitos de información}
Para modelar fielmente la red logística, el sistema debe ser capaz de persistir y relacionar entidades que poseen atributos tanto técnicos como descriptivos para su visualización:

\begin{itemize}
    \item \textbf{Almacenes (Nodos):} Además de su identificación y capacidad, deben incluir metadatos de localización (como ciudad o coordenadas) que permitan su representación en el mapa interactivo.
    \item \textbf{Productos (Nodos):} Deben categorizarse mediante identificadores únicos (SKU), nombres descriptivos y categorías para facilitar su filtrado en el inventario.
    \item \textbf{Conexiones Viales (Aristas):} Deben definir la topología de la red, vinculando almacenes mediante el atributo de distancia, esencial para los algoritmos de optimización.
    \item \textbf{Existencias (Aristas):} Representan la relación dinámica entre productos y almacenes, manteniendo un conteo preciso de unidades disponibles en cada punto de la red.
\end{itemize}

\section{Requisitos funcionales}
Hemos definido las capacidades del sistema centrándonos en la utilidad para el administrador final, equilibrando la potencia de cálculo con la facilidad de uso:

\begin{itemize}
    \item \textbf{Gestión de Red Interactiva:} El sistema debe permitir realizar operaciones CRUD sobre almacenes, productos y rutas. Estas acciones deben reflejarse inmediatamente tanto en la base de datos como en la representación visual de la red.
    \item \textbf{Cálculo de Rutas Críticas:} Implementar una lógica capaz de determinar el camino más corto entre dos puntos, garantizando que el algoritmo recorra exclusivamente las rutas de transporte válidas.
    \item \textbf{Dashboard de Inventario:} Proporcionar una visión clara del stock vivo, permitiendo realizar ajustes de unidades y visualizar la distribución de productos a través de la red.
    \item \textbf{Visualización Dinámica de la Red:} La interfaz debe generar un grafo interactivo que ayude al usuario a comprender la topología de la red logística de un solo vistazo.
    \item \textbf{Resiliencia ante Contingencias:} Capacidad de recalcular rutas de forma automática si una conexión queda inhabilitada, asegurando la continuidad de la cadena de suministro.
\end{itemize}

\section{Requisitos no funcionales}
Para asegurar la calidad técnica y la facilidad de evaluación del proyecto, hemos establecido los siguientes criterios:
\begin{itemize}
    \item \textbf{Portabilidad:} El sistema debe ser agnóstico al entorno de ejecución, garantizando su funcionamiento inmediato mediante el uso de contenedores.
    \item \textbf{Rendimiento en Relaciones:} La arquitectura debe permitir realizar consultas de vecindad y caminos cortos sin la penalización de rendimiento propia de los sistemas relacionales convencionales.
    \item \textbf{Persistencia y Consistencia:} Se debe asegurar que las operaciones sobre el grafo mantengan la integridad de la red, evitando nodos o aristas huérfanas tras una eliminación.
\end{itemize}
